% Part: incompleteness
% Chapter: representability-in-q
% Section: comp-representable

\documentclass[../../../include/open-logic-section]{subfiles}

\begin{document}

\olfileid{inc}{req}{crq}
\olsection{محاسبه کېدونکې تابعې په~$\Th{Q}$ کښې تمثيلېدونکې دي}

\begin{thm}
هره محاسبه کېدونکې تابعه په~$\Th{Q}$ کښې تمثيلېدونکې ده۔
\end{thm}

\begin{proof}
د کره والي دپاره، او د چرچ--ټيورينګ د دعوې په کارولو سره، ووايو چې
يوه تابعه هله او يوازې هله محاسبه کېدونکې ده چې عمومي بازګشتي وي۔
عمومي بازګشتي تابعې هغه دي چې له صفر تابع~$\Zero$، تالي
تابع~$\Succ$، او پروجکشن تابع~$\Proj{n}{i}$ څخه د ترکيب، بنسټيز
بازګښت، او منظم اقل موندلو په کارولو تعريفېدے شي۔ د
\olref[pri]{lem:prim-rec} له مخې، هره تابعه~$h$ چې له~$f$ او~$g$
څخه تعريفېدے شي، له~$f$، $g$، او~$\Zero$، $\Succ$،
$\Proj{n}{i}$، $\Add$، $\Mult$، $\Char{=}$ څخه هم د ترکيب او منظم
اقل موندلو په کارولو تعريفېدے شي۔ له دې امله، يوه تابعه هله او يوازې
هله عمومي بازګشتي ده چې له~$\Zero$، $\Succ$، $\Proj{n}{i}$،
$\Add$، $\Mult$، $\Char{=}$ څخه د ترکيب او منظم اقل موندلو په
کارولو تعريفېدے شي۔

سربېره پر دې، ښودلې مو ده چې د پام وړ بنسټيزې تابعې په~$\Th{Q}$
کښې تمثيلېدونکې دي
(\cref{inc:req:bre:prop:rep-zero,inc:req:bre:prop:rep-succ,inc:req:bre:prop:rep-proj,inc:req:bre:prop:rep-id,inc:req:bre:prop:rep-add,inc:req:bre:prop:rep-mult})،
او هره هغه تابعه چې له تمثيلېدونکو تابعو څخه په ترکيب يا منظم اقل
موندلو تعريف شوې وي
(\olref[cmp]{prop:rep-composition}،
\olref[min]{prop:rep-minimization}) هم تمثيلېدونکې ده۔ نو هره عمومي
بازګشتي تابعه په~$\Th{Q}$ کښې تمثيلېدونکې ده۔
\end{proof}

\begin{explain}
موږ ښودلې ده چې د محاسبه کېدونکو تابعو سټ د هغو تابعو د سټ په توګه
ځانګړے کېدے شي چې په~$\Th{Q}$ کښې تمثيلېدونکې دي۔ په حقيقت کښې
ثبوت تر دې عام دے۔ د تمثيلېدنې له تعريف څخه په اسانۍ ليدل کېږي چې
هره هغه تيوري چې~$\Th{Q}$ غځوي (يا~$\Th{Q}$ پکې تفسيرېدے شي)،
محاسبه کېدونکې تابعې تمثيلولے شي۔ خو برعکس، په هر هغه
!!{derivation} نظام کښې چې د !!{derivation} مفهوم ئې محاسبه کېدونکے
وي، هره تمثيلېدونکې تابعه محاسبه کېدونکې ده۔ نو د بېلګې په توګه، د
محاسبه کېدونکو تابعو سټ د هغو تابعو د سټ په توګه ځانګړے کېدے شي چې
د پيانو په حساب، يا ان د زرمیلو--فرېنکل په سټ تيورۍ کښې تمثيلېدونکې
دي۔ لکه ګوډل چې يادونه کړې، دا يو څه هېښوونکې ده۔ وبه وينو چې د
ثبوت‌وړتيا په برخه کښې پوښتنې د ټاکلې تيورۍ پر وړاندې ډېرې حساسې
دي؛ په ټوليز ډول هر څومره چې بديهي اصول پياوړي وي، هغومره ډېر څه
ثابتولے شئ۔ خو د بديهي تيوريو په يوه پراخه ډله کښې تمثيلېدونکې
تابعې کټ مټ محاسبه کېدونکې تابعې دي؛ تر هغه چې تيورۍ
بديهي کېدونکې وي، پياوړې تيورۍ زياتې تابعې نۀ تمثيلوي۔
\end{explain}

\end{document}
